\vspace{-2mm}\section{Types of Reasoning Hallucination}\vspace{-1mm}
\label{sec:taxonomy}

Based on the formalization in Definition~\ref{def:reasoning_hallucination}, we systematically categorize reasoning hallucination ($H_R$) into two primary classes.

\subsection{Invalidity Flaws}
\label{sec:invalidity_flaws}

Invalidity Flaws are violations of Valid Reasoning (Definition~\ref{def:valid_reasoning}).
These flaws manifest when at least one reasoning step $r_i$ is logically invalid (i.e., $V(r_i | r_{<i}, Q) = \text{False}$), or when the complete reasoning process $R$ fails to logically entail the final answer $A$ (i.e., $C(R, A) = \text{False}$).

This category encompasses several distinct failure modes.
\textbf{Logical Fallacies} represent the most direct form of invalidity ($V(r_i) = \text{False}$), where models produce steps that are not logically sound.
This includes causal misattributions, such as the ``illusion of causality'' or the ``post hoc fallacy''~\cite{joshi2024llms}, and a ``Fake Reasoning Bias''~\cite{wang2025towards}, where models favor superficial rhetorical structures over logical soundness.

Beyond single-step failures, \textbf{Think-Answer Mismatch} corresponds to the $C(R, A) = \text{False}$ condition.
In this phenomenon, the model's generated reasoning chain $R$ is disconnected from its final answer $A$, suggesting the CoT serves as a post-hoc rationalization rather than the true derivation process~\cite{yao2025reasoning}.
Notably, even when a model exhibits this type of hallucination and reaches an incorrect conclusion during the reasoning process, it may still produce the correct final answer.
This possibility poses a significant challenge to the identification of reasoning hallucination~\cite{yee2024dissociation}.

A distinct type of invalidity emerges as \textbf{Flaw Repetition}, where the model becomes trapped in a loop of semantically similar but incorrect reasoning steps~\cite{yao2025reasoning}.

\subsection{Sub-optimality Flaws}
\label{sec:suboptimality_flaws}

Sub-optimality Flaws represent a more insidious form of $H_R$.
As defined in Definition~\ref{def:reasoning_hallucination}, these are processes that are technically $Valid$ but are pathological due to redundancy (i.e., $\exists R' \subset R \text{ s.t. } Valid(R', Q, A)$).

The primary manifestation of this category is \textbf{Overthinking}.
This phenomenon is characterized by the generation of verbose and tangential reasoning traces even for simple queries~\cite{cuesta2025lrms}.
More pathologically, this behavior causes models to disregard correct solutions even when explicitly provided, leading them to pursue unnecessarily complex and often erroneous reasoning paths.
This suggests that the model is pathologically matching the form of complex reasoning (i.e., generating long token sequences) rather than executing the function of efficient problem-solving.
